\section{Observabilidad y control de seguridad del AI Agent}

\subsection{Indicadores de monitoreo y estrategia de alertas}

Una plataforma de Agent madura primero debe definir indicadores observables y después generar alertas automáticas en cuanto aparezcan «señales»:

\begin{itemize}
    \item \textbf{Consistencia semántica}: Historical Reference vs. current response divergence
    \item \textbf{Tasa de alucinación del modelo}: revisar aleatoriamente las salidas para ver si se introdujeron hechos falsos
    \item \textbf{Errores en la llamada a herramientas}: fallas en la ejecución de comandos / tiempos de espera agotados en llamadas a la API
    \item \textbf{Explosión del contexto}: una expansión repentina de la ventana de contexto indica un memory leak o una prompt injection
    \item \textbf{Número de correcciones del usuario}: interrupciones frecuentes del usuario indican que el agent malinterpretó el intent
\end{itemize}

\begin{knowledgebox}{Cuanto antes se capturen las signals, más fácil es corregir el rumbo}
Muchos equipos solo se enteran de que el agent tiene un problema cuando el usuario se queja. Las señales anómalas tempranas suelen esconderse en un latency spike, en tiempos de espera agotados en la llamada a herramientas o en «contenido repetido». Establecer un esquema transparente de «metric + sampling reforzado + human-in-loop review» permite escalar la respuesta antes de que el agent falle.
\end{knowledgebox}

Burak subraya que el pipeline de datos de instrumentation debe mantener un costo operativo bajo: para la consistencia semántica o la hallucination, solo se muestrea entre el 1 y el 2\% de las solicitudes, y la sensibilidad se preserva mediante la revisión del squad. Recolectar en exceso genera costos y también puede provocar fatiga de alertas.

\subsection{Gobernanza de seguridad y diseño de guardrail}

El comportamiento del Agent no puede depender únicamente del prompt del modelo; también depende de guardrails a nivel de sistema:

\begin{itemize}
    \item \textbf{Policy Layer}: intercepta operaciones sensibles antes de ejecutarlas (eliminar archivos, programar tareas)
    \item \textbf{Tool Firewall}: limita el acceso del modelo a las herramientas a una lista blanca
    \item \textbf{Human-in-the-loop level}: decide dinámicamente, según el riesgo de la tarea, si se necesita aprobación humana
    \item \textbf{Audit Trail}: registra cada decisión, cada llamada a herramientas y cada confirmación del usuario
\end{itemize}

\begin{warningbox}{Sin audit trail no hay manera de exigir responsabilidades}
En cuanto un agent provoca un error o un incidente de seguridad, la falta de un registro de operaciones significa que no se puede reconstruir la cadena de decisiones. Esto hace que el sistema de gobernanza deje de funcionar e incluso afecta la auditoría de cumplimiento. Tanto Anthropic como Vertex AI subrayan que cada tool call debe ser trazable y reversible.
\end{warningbox}

\subsection{Resumen de la sección}

Solo mediante indicadores cuantitativos, cadenas de alerta y cadenas de responsabilidad puede una empresa reaccionar con rapidez cuando el agent presenta anomalías, y mantener así los límites de seguridad al mismo tiempo que crece su capacidad.

